\chapter{Gröbner-Basen}

\section{Die Basis, die alles sieht}

\begin{merksatz}
\boxtitel{Der Leitfall in drei Zeilen}
Aus $f_1, f_2, f_3$ berechnet der Algorithmus dieses Kapitels in lex-Ordnung ($k_{10} > k_{12} > k_{21} > k_{20}$):
\[ G = \bigl\{\ k_{10} + k_{12} - \tfrac12,\quad k_{21} + k_{20} - \tfrac25,\quad k_{12}k_{20} - \tfrac25k_{12} + \tfrac{3}{50}\ \bigr\} . \]
\begin{center}
{\small
\renewcommand{\arraystretch}{1.5}
\begin{tabular}{@{}P{3.6cm}P{5.0cm}@{}}
\toprule
\textbf{Leitmonome} & $k_{10}$, $k_{21}$, $k_{12}k_{20}$\\
\textbf{Was sie sagen} & $k_{10}$ und $k_{21}$ sind durch die anderen bestimmt; zwischen $k_{12}$ und $k_{20}$ gilt eine Gleichung; ein Parameter bleibt frei.\\
\textbf{Zugehörigkeit} & $f\in I$ genau dann, wenn Rest bei Division durch $G$ null.\\
\textbf{Dimension} & Freie Veränderliche: eine. Also eine Kurve.\\
\bottomrule
\end{tabular}}
\end{center}

\begin{center}
\begin{minipage}{0.9\textwidth}
\textbf{Die Gröbner-Basis ist für Polynomsysteme, was die Zeilenstufenform für lineare Systeme ist: eine Normalform, an der man Lösbarkeit, Dimension und Freiheitsgrade abliest --- und die der Rechner aus jedem Erzeugendensystem berechnet.}
\end{minipage}
\end{center}

\begin{center}
\begin{minipage}{0.88\textwidth}
\kursiv{\textbf{Drei Dinge beweist dieses Kapitel}: dass eine solche Basis für jedes Ideal existiert (schon in Kapitel 2 geschehen), dass man sie an einem Kriterium erkennt, das nur endlich viele Rechnungen braucht (Buchberger-Kriterium), und dass ein Algorithmus sie findet und endet (Buchberger-Algorithmus). Dazu, dass es eine kanonische Form gibt --- die reduzierte Basis ---, die vom Erzeugendensystem nicht mehr abhängt. Damit wird aus der Frage \enquote{sind zwei Ideale gleich?} ein Vergleich zweier Listen.}
\end{minipage}
\end{center}
\end{merksatz}

\section{Definition und erste Sätze}

\begin{definition}
\boxtitelm{Definition 5.1 --- Gröbner-Basis}
Sei eine Monomordnung fixiert und $I\ne\{0\}$ ein Ideal. Eine endliche Teilmenge $G = \{g_1, \ldots, g_s\}\subseteq I$ heißt \Begriff{Gröbner-Basis} von $I$, wenn
\[ \langle\mathrm{LM}(g_1), \ldots, \mathrm{LM}(g_s)\rangle = \mathrm{LM}(I) := \langle\mathrm{LM}(f) : f\in I\rangle . \]
Sie heißt \Begriff{reduziert}, wenn alle $\mathrm{LC}(g_i) = 1$ und kein Monom eines $g_i$ durch ein $\mathrm{LM}(g_j)$, $j\ne i$, teilbar ist.
\end{definition}

\begin{merksatz}
\boxtitel{Satz 5.2 --- was eine Gröbner-Basis leistet}
Sei $G$ eine Gröbner-Basis von $I$.
\begin{enumerate}[label=(\alph*), leftmargin=1.6em]
\item $G$ erzeugt $I$.
\item Der Rest bei Division durch $G$ ist eindeutig --- unabhängig von der Reihenfolge der $g_i$. Man schreibt $\overline f^{\,G}$.
\item $f\in I$ genau dann, wenn $\overline f^{\,G} = 0$.
\item Jedes Ideal hat genau eine reduzierte Gröbner-Basis (zu fester Ordnung).
\end{enumerate}

\schritt{Beweis von (a)}
Das ist Schritt 3 im Beweis des Basissatzes (Satz 2.3): Für $f\in I$ hat der Rest $r$ bei Division durch $G$ kein Monom in $\mathrm{LM}(I)$, liegt aber in $I$; also $r = 0$. $\blacksquare$

\schritt{Beweis von (b)}
Seien $f = \sum q_ig_i + r = \sum q_i'g_i + r'$ zwei Divisionen. Dann ist $r - r'\in I$, und kein Monom von $r - r'$ ist durch ein $\mathrm{LM}(g_i)$ teilbar. Wäre $r - r'\ne0$, so läge $\mathrm{LM}(r - r')\in\mathrm{LM}(I) = \langle\mathrm{LM}(g_i)\rangle$, wäre also durch ein $\mathrm{LM}(g_i)$ teilbar. Widerspruch; $r = r'$. $\blacksquare$

\schritt{Beweis von (c)}
$\Leftarrow$: $f = \sum q_ig_i\in I$. $\Rightarrow$: wie (a). $\blacksquare$

\schritt{Beweis von (d), Existenz}
Man entferne aus $G$ jedes $g_i$, dessen Leitmonom durch ein anderes $\mathrm{LM}(g_j)$ teilbar ist (die übrigen Leitmonome erzeugen immer noch $\mathrm{LM}(I)$). Dann ersetze man jedes $g_i$ durch $\overline{g_i}^{\,G\setminus\{g_i\}}$ --- den Rest bei Division durch die anderen --- und normiere. Das Leitmonom bleibt (es ist durch kein anderes teilbar, landet also im Rest als erstes), alle anderen Monome werden reduziert. $\blacksquare$

\schritt{Beweis von (d), Eindeutigkeit}
Seien $G, G'$ reduzierte Gröbner-Basen. Die Leitmonome beider sind die minimalen Erzeuger des Monomideals $\mathrm{LM}(I)$, und die sind eindeutig: ein Monomideal hat genau eine minimale Erzeugermenge (die Monome, die durch kein anderes Element teilbar sind). Also gibt es zu jedem $g\in G$ ein $g'\in G'$ mit $\mathrm{LM}(g) = \mathrm{LM}(g')$. Dann ist $g - g'\in I$, und $g - g'$ enthält kein Monom aus $\mathrm{LM}(I)$ (die Leitterme heben sich weg, die übrigen Monome sind nach Reduziertheit durch kein Leitmonom teilbar). Nach (c) und (b) ist $\overline{g - g'}^{\,G} = g - g' = 0$. $\blacksquare$

\begin{center}
\begin{minipage}{0.88\textwidth}
\kursiv{\textbf{(d) ist der Satz, der aus Idealen Daten macht}: Zwei Erzeugendensysteme erzeugen genau dann dasselbe Ideal, wenn ihre reduzierten Gröbner-Basen übereinstimmen --- ein Zeichenvergleich. Im Leitfall erzeugen $\{f_1, f_2, f_3\}$ und $\{f_1, f_2 - f_1, g\}$ dasselbe Ideal; beide haben die reduzierte Basis $G$ oben. Und der eindeutige Rest $\overline f^{\,G}$ ist die Normalform von $f$ modulo $I$: Der Koordinatenring $k[V]$ (Kapitel 3) wird zum Vektorraum der Reste, mit den Monomen außerhalb von $\mathrm{LM}(I)$ als Basis --- den \Begriff{Standardmonomen}. Kapitel 7 zählt sie.}
\end{minipage}
\end{center}
\end{merksatz}

\section{Das Kriterium}

\begin{definition}
\boxtitelm{Definition 5.3 --- S-Polynom}
Für $f, g\ne0$ sei $x^{\gamma} = \mathrm{kgV}(\mathrm{LM}(f), \mathrm{LM}(g))$. Das \Begriff{S-Polynom} ist
\[ S(f, g) = \frac{x^{\gamma}}{\mathrm{LT}(f)}\,f - \frac{x^{\gamma}}{\mathrm{LT}(g)}\,g . \]
Es ist so gebaut, dass die Leitterme sich wegheben: $\mathrm{LM}(S(f, g)) < x^{\gamma}$.
\end{definition}

\begin{merksatz}
\boxtitel{Satz 5.4 --- Buchbergers Kriterium}
$G = \{g_1, \ldots, g_s\}$ ist genau dann eine Gröbner-Basis von $\langle G\rangle$, wenn für alle $i\ne j$
\[ \overline{S(g_i, g_j)}^{\,G} = 0 . \]

\schritt{Beweis: $\Rightarrow$}
$S(g_i, g_j)\in I$, also Rest null nach Satz 5.2(c). $\blacksquare$

\schritt{Beweis: $\Leftarrow$, Vorbereitung}
Sei $f\in I$, $f\ne0$; zu zeigen: $\mathrm{LM}(f)\in\langle\mathrm{LM}(g_i)\rangle$. Schreibe $f = \sum h_ig_i$ und setze $\delta = \max_i\mathrm{LM}(h_ig_i)$. Unter allen Darstellungen wähle eine mit kleinstem $\delta$ (Wohlordnung). Es ist $\mathrm{LM}(f)\le\delta$. Ist $\mathrm{LM}(f) = \delta$, so ist $\mathrm{LM}(f) = \mathrm{LM}(h_ig_i)$ für ein $i$, also durch $\mathrm{LM}(g_i)$ teilbar --- fertig. Bleibt der Fall $\mathrm{LM}(f) < \delta$: Die Terme vom Monom $\delta$ löschen sich aus. Wir zeigen, dass dann eine Darstellung mit kleinerem $\delta$ existiert --- Widerspruch zur Minimalität.

\schritt{Beweis: $\Leftarrow$, der Kern}
Seien $i_1, \ldots, i_m$ die Indizes mit $\mathrm{LM}(h_ig_i) = \delta$. Setze $c_i = \mathrm{LC}(h_i)$ und $x^{\alpha(i)} = \mathrm{LM}(h_i)$, also $x^{\alpha(i)}\mathrm{LM}(g_i) = \delta$. Dann ist $\sum_{i\in\{i_1..i_m\}}c_i\mathrm{LC}(g_i) = 0$ (Auslöschung). Lemma: Eine Summe $\sum c_ix^{\alpha(i)}g_i$ mit $x^{\alpha(i)}\mathrm{LM}(g_i) = \delta$ für alle $i$ und $\sum c_i\mathrm{LC}(g_i) = 0$ ist eine Linearkombination $\sum_{j<l}c_{jl}\,x^{\delta - \gamma_{jl}}S(g_j, g_l)$ mit $x^{\gamma_{jl}} = \mathrm{kgV}(\mathrm{LM}(g_j), \mathrm{LM}(g_l))$. Beweis des Lemmas: Setze $p_i = x^{\alpha(i)}g_i/\mathrm{LC}(g_i)$ (Leitterm $\delta$, Koeffizient $1$) und $d_i = c_i\mathrm{LC}(g_i)$, $\sum d_i = 0$. Dann $\sum c_ix^{\alpha(i)}g_i = \sum d_ip_i = d_{i_1}(p_{i_1} - p_{i_2}) + (d_{i_1} + d_{i_2})(p_{i_2} - p_{i_3}) + \cdots$ (Teleskopsumme; der letzte Koeffizient ist $\sum d_i = 0$). Und $p_j - p_l = x^{\delta - \gamma_{jl}}S(g_j, g_l)$, weil beide Seiten dieselbe Kombination mit gleichem Leitterm sind (nachrechnen: $x^{\alpha(j)}/\mathrm{LC}(g_j) = x^{\delta - \gamma_{jl}}\cdot x^{\gamma_{jl}}/\mathrm{LT}(g_j)$). $\square$

\schritt{Beweis: $\Leftarrow$, Abschluss}
Nach Voraussetzung ist $\overline{S(g_j, g_l)}^{\,G} = 0$, also $S(g_j, g_l) = \sum_i a_{jli}g_i$ mit $\mathrm{LM}(a_{jli}g_i)\le\mathrm{LM}(S(g_j, g_l)) < x^{\gamma_{jl}}$ (Satz 4.3). Multipliziert mit $x^{\delta - \gamma_{jl}}$: alle Terme $< \delta$. Ersetzt man in der Darstellung von $f$ die Terme vom Monom $\delta$ durch diese Ausdrücke, erhält man eine Darstellung $f = \sum h_i'g_i$, in der alle $\mathrm{LM}(h_i'g_i) < \delta$. Widerspruch zur Minimalität von $\delta$. $\blacksquare$

\begin{center}
\begin{minipage}{0.9\textwidth}
\textbf{Der Beweis sagt, warum S-Polynome genügen: Jede Auslöschung von Leittermen in einer Summe ist eine Kombination von paarweisen Auslöschungen --- das Lemma. Wer die Paare beherrscht, beherrscht alle Summen.}
\end{minipage}
\end{center}
\end{merksatz}

\section{Der Algorithmus}

\begin{merksatz}
\boxtitel{Satz 5.5 --- Buchbergers Algorithmus}
Eingabe $F = \{f_1, \ldots, f_s\}$. Setze $G = F$. Solange es ein Paar $g_i, g_j\in G$ gibt mit $r = \overline{S(g_i, g_j)}^{\,G}\ne0$: füge $r$ zu $G$ hinzu. Ausgabe $G$. Der Algorithmus endet, und $G$ ist eine Gröbner-Basis von $\langle F\rangle$.

\schritt{Korrektheit}
Jedes hinzugefügte $r$ liegt in $\langle F\rangle$ (es ist eine Kombination von Elementen aus $G$), also bleibt $\langle G\rangle = \langle F\rangle$. Beim Abbruch haben alle S-Polynome Rest null: Gröbner-Basis nach Satz 5.4. $\blacksquare$

\schritt{Ende}
Jedes hinzugefügte $r$ hat ein Leitmonom, das durch kein bisheriges $\mathrm{LM}(g_i)$ teilbar ist (sonst wäre es nicht im Rest). Also wächst das Monomideal $\langle\mathrm{LM}(G)\rangle$ bei jedem Schritt echt. Eine echt aufsteigende Kette von Monomidealen ist endlich (Dickson / Basissatz). $\blacksquare$

\schritt{Zwei Abkürzungen}
\begin{center}
{\small
\renewcommand{\arraystretch}{1.5}
\begin{tabular}{@{}P{3.0cm}P{5.8cm}@{}}
\toprule
\textbf{Teilerfremde Leitmonome} & Ist $\mathrm{kgV}(\mathrm{LM}(g_i), \mathrm{LM}(g_j)) = \mathrm{LM}(g_i)\mathrm{LM}(g_j)$, so ist $\overline{S(g_i, g_j)}^{\,G} = 0$ von selbst (Buchbergers erstes Kriterium) --- das Paar muss nicht gerechnet werden.\\
\textbf{Kettenkriterium} & Ist $\mathrm{LM}(g_k)$ ein Teiler von $\mathrm{kgV}(\mathrm{LM}(g_i), \mathrm{LM}(g_j))$ und sind die Paare $(i,k)$, $(k,j)$ erledigt, so ist $(i,j)$ überflüssig (zweites Kriterium).\\
\bottomrule
\end{tabular}}
\end{center}

\schritt{Die verdrehte Kubik, Schritt für Schritt (lex, $x > y > z$)}
\begin{center}
{\small
\renewcommand{\arraystretch}{1.46}
\begin{tabular}{@{}P{1.4cm}P{3.0cm}P{4.2cm}@{}}
\toprule
\textbf{Paar} & $S$ & \textbf{Rest durch $G$}\\
\midrule
$(g_1, g_2)$ & $x\,g_1 - g_2 = -xy + z$ & $-xy + z$ --- neu: $g_3 = xy - z$\\
$(g_1, g_3)$ & $y\,g_1 - x\,g_3 = -y^{2} + xz$ & $xz - y^{2}$ --- neu: $g_4 = xz - y^{2}$\\
$(g_2, g_3)$ & $y\,g_2 - x^{2}g_3 = x^{2}z - yz$ & $x^{2}z - yz\to z\,g_1 = x^{2}z - yz$: Rest $0$\\
$(g_1, g_4)$ & $z\,g_1 - x\,g_4 = xy^{2} - yz$ & $y\,g_3 = xy^{2} - yz$: Rest $0$\\
$(g_2, g_4)$ & $z\,g_2 - x^{2}g_4 = x^{2}y^{2} - z^{2}$ & $y^{2}g_1 = x^{2}y^{2} - y^{3}$; Rest $y^{3} - z^{2}$ --- neu: $g_5$\\
$(g_3, g_4)$ & $z\,g_3 - y\,g_4 = y^{3} - z^{2}$ & $= g_5$: Rest $0$\\
übrige & teilerfremd oder Kette & Rest $0$\\
\bottomrule
\end{tabular}}
\end{center}
Reduziert ($g_2 = x\,g_1 + g_3$ fliegt heraus): $G = \{x^{2} - y,\ xy - z,\ xz - y^{2},\ y^{3} - z^{2}\}$. Das ist die Basis aus Kapitel 2. \checkmark

\begin{center}
\begin{minipage}{0.88\textwidth}
\kursiv{\textbf{Das Programm, mit dem dieser Band rechnet, ist genau Satz 5.5 mit dem ersten Kriterium}: hundert Zeilen --- Polynome als Wörterbücher von Exponentenvektoren, eine Vergleichsfunktion für die Ordnung, Division, S-Polynom, Schleife, Reduktion. Es rechnet alle Beispiele des Bandes in Sekundenbruchteilen und ist für mehr nicht gedacht; die Systeme, die man wirklich benutzt (Singular, Macaulay2, Magma, die Bibliotheken in Sage und Maple), haben dieselbe Struktur und zwanzig Jahre Optimierung: Paarauswahl nach Grad, lineare Algebra über viele S-Polynome auf einmal (F4, \Person{Faugère} 1999), Signaturen, die überflüssige Reduktionen vorab erkennen (F5, 2002). Der Leser, der das Programm nach diesem Kapitel schreibt, hat den Algorithmus verstanden; wer es schnell machen will, hat ein Forschungsgebiet vor sich.}
\end{minipage}
\end{center}
\end{merksatz}

\section{Der Bruchfall}

\begin{warnung}
\boxtitelr{Die Basis hängt von der Ordnung ab --- und kann riesig werden}
Dasselbe Ideal $\langle x^{2} + y^{2} + z^{2} - 4,\ x^{2} + 2y^{2} - 5,\ xz - 1\rangle$:
\begin{center}
{\small
\renewcommand{\arraystretch}{1.5}
\begin{tabular}{@{}P{1.4cm}P{7.4cm}@{}}
\toprule
grevlex & $\{z^{3} + \tfrac12x - \tfrac32z,\ x^{2} + 2z^{2} - 3,\ y^{2} - z^{2} - 1,\ xz - 1\}$\\
lex & $\{x + 2z^{3} - 3z,\ y^{2} - z^{2} - 1,\ z^{4} - \tfrac32z^{2} + \tfrac12\}$\\
\bottomrule
\end{tabular}}
\end{center}
Die lex-Basis ist dreieckig und löst das System: $z^{2} = 1$ oder $\tfrac12$, dann $y$, dann $x$ --- acht Lösungen. Die grevlex-Basis löst nichts, ist aber kleiner im Grad.

\schritt{Und für vier Veränderliche}
Das Ideal \enquote{cyclic-4} ($a + b + c + d$, $ab + bc + cd + da$, $abc + bcd + cda + dab$, $abcd - 1$): grevlex-Basis sieben Elemente, höchster Grad $6$; lex-Basis sechs Elemente, höchster Grad $8$, darunter $c^{2}d^{6} - c^{2}d^{2} - d^{4} + 1$. Für cyclic-7 hat die lex-Basis Elemente vom Grad über $200$ und Koeffizienten mit Hunderten Stellen; cyclic-9 wurde erst 2001 bewältigt.

\begin{center}
\begin{minipage}{0.9\textwidth}
\textbf{Die Gröbner-Basis ist eine Normalform, aber ihre Größe ist keine Eigenschaft des Ideals allein, sondern des Paars (Ideal, Ordnung). Im schlimmsten Fall ist sie doppelt exponentiell in der Zahl der Veränderlichen --- Mayr--Meyer, Kapitel 2 ---, und lex ist regelmäßig der schlimmere Fall.}
\end{minipage}
\end{center}

\begin{center}
\begin{minipage}{0.88\textwidth}
\kursiv{\textbf{Der Leitfall ist gutartig --- zwei lineare Gleichungen, eine quadratische ---, und das ist typisch für Kompartimentmodelle, nicht für Polynomsysteme im Allgemeinen.} Die Praxis zieht drei Lehren: grevlex rechnen und bei Bedarf umordnen (Kapitel 4, Forschungskasten); Symmetrie und Struktur vor dem Rechnen ausnutzen (cyclic-$n$ ist symmetrisch, und wer das benutzt, rechnet cyclic-10); und den Grad überwachen --- ein Buchberger-Lauf, dessen Zwischenpolynome den Grad $50$ überschreiten, endet meist nicht in nützlicher Zeit. Dass derselbe Algorithmus, der in Kapitel 1 eine Hyperbel fand, in anderen Fällen hoffnungslos ist, gehört zur Wahrheit über ihn.}
\end{minipage}
\end{center}
\end{warnung}

\section{Übungsaufgaben}

\begin{aufgabe}
\boxtitel{Aufgabe 5.1 --- Eine Basis von Hand}
Berechnen Sie die reduzierte Gröbner-Basis von $\langle xy - 1,\ y^{2} - 1\rangle$ in lex, $x > y$ (die Erzeuger aus Aufgabe 4.2).

\loesung
\schritt{Schritt 1 --- S-Polynom}
$\mathrm{kgV}(xy, y^{2}) = xy^{2}$: $S = y(xy - 1) - x(y^{2} - 1) = -y + x = x - y$.

\schritt{Schritt 2 --- Rest}
$\mathrm{LM}(x - y) = x$, durch $xy$ und $y^{2}$ nicht teilbar: Rest $x - y\ne0$. Neu: $g_3 = x - y$.

\schritt{Schritt 3 --- Neue Paare}
$S(g_1, g_3) = (xy - 1) - y(x - y) = y^{2} - 1 = g_2$: Rest $0$. $S(g_2, g_3)$: Leitmonome $y^{2}$ und $x$ teilerfremd --- erstes Kriterium, Rest $0$.

\schritt{Schritt 4 --- Reduzieren}
$\mathrm{LM}(g_1) = xy$ ist durch $\mathrm{LM}(g_3) = x$ teilbar: $g_1$ fliegt heraus. $g_2 = y^{2} - 1$ und $g_3 = x - y$ sind gegeneinander reduziert. $G = \{x - y,\ y^{2} - 1\}$.

\schritt{Schritt 5 --- Probe}
$V(G)$: $y = \pm1$, $x = y$: die Punkte $(1,1)$ und $(-1,-1)$. Im Original: $xy = 1$ und $y^{2} = 1$ geben dieselben. \checkmark\ Der Rest von $f = x^{2}y + xy^{2} + y^{2}$ durch $G$: $x^{2}y\to$ (ersetze $x$ durch $y$) $y^{3}\to y$; $xy^{2}\to y^{3}\to y$; $y^{2}\to1$: Rest $2y + 1$ --- der kanonische Rest aus Aufgabe 4.2. \checkmark
\end{aufgabe}

\begin{aufgabe}
\boxtitel{Aufgabe 5.2 --- Sind zwei Ideale gleich?}
$I = \langle x^{2} + y,\ x + y^{2}\rangle$ und $J = \langle x + y^{2},\ y^{4} + y\rangle$ in $\mathbb Q[x, y]$, lex $x > y$. Gleich?

\loesung
\schritt{Schritt 1 --- Basis von $I$}
$S(x^{2} + y, x + y^{2}) = (x^{2} + y) - x(x + y^{2}) = -xy^{2} + y$. Rest durch $\{x^{2} + y, x + y^{2}\}$: $-xy^{2}$ ist durch $x$ teilbar: $-xy^{2} + y + y^{2}(x + y^{2}) = y^{4} + y$; kein Leitmonom passt mehr: Rest $y^{4} + y$. Neu. Weitere S-Polynome: $S(x + y^{2}, y^{4} + y)$ teilerfremd; $S(x^{2} + y, y^{4} + y) = y^{4}(x^{2} + y) - x^{2}(y^{4} + y) = y^{5} - x^{2}y\to$ Rest: $-x^{2}y + y^{5} + y(x^{2} + y) = y^{5} + y^{2} = y(y^{4} + y)$: Rest $0$. Reduzieren: $x^{2} + y$ fliegt heraus ($x^{2}$ durch $x$ teilbar). Reduzierte Basis $\{x + y^{2},\ y^{4} + y\}$.

\schritt{Schritt 2 --- Basis von $J$}
Die Erzeuger $x + y^{2}$, $y^{4} + y$ haben teilerfremde Leitmonome $x$, $y^{4}$ und sind gegeneinander reduziert: Das ist schon die reduzierte Basis.

\schritt{Schritt 3 --- Vergleich}
Beide reduzierten Basen sind $\{x + y^{2}, y^{4} + y\}$: $I = J$ (Satz 5.2d). \checkmark

\schritt{Schritt 4 --- Geometrie}
$y^{4} + y = y(y^{3} + 1) = y(y + 1)(y^{2} - y + 1)$: $y\in\{0, -1, \tfrac12\pm\tfrac{\sqrt3}{2}\mathrm i\}$, $x = -y^{2}$. Vier Punkte, zwei reell: $(0,0)$ und $(-1,-1)$. Parabel $y = -x^{2}$ und Parabel $x = -y^{2}$ schneiden sich in vier Punkten --- Bézout, Kapitel 9, $2\cdot2 = 4$. \checkmark

\begin{center}
\begin{minipage}{0.88\textwidth}
\kursiv{\textbf{Ohne Satz 5.2(d) wäre die Frage \enquote{$I = J$?} eine Suche nach zwei Zertifikaten in jede Richtung}; mit ihm ist sie ein Vergleich. Das ist der praktische Kern des Kapitels, und er trägt bis in die Statistik: Zwei Modelle sind algebraisch äquivalent, wenn ihre Ideale gleich sind, Kapitel 13.}
\end{minipage}
\end{center}
\end{aufgabe}

\begin{aufgabe}
\boxtitel{Aufgabe 5.3 --- Der Leitfall mit Messung im Gewebe}
Für das Modell aus Aufgabe 1.3 (Messung in Kammer 2, $k_{20} = 0$) mit Daten $k_{12} = \tfrac15$, $k_{10} + k_{12} + k_{21} = \tfrac45$, $k_{10}k_{21} = \tfrac{8}{100}$: Berechnen Sie die reduzierte Gröbner-Basis (lex, $k_{10} > k_{12} > k_{21}$) und lesen Sie die Identifizierbarkeit ab.

\loesung
\schritt{Schritt 1 --- Erzeuger}
$h_1 = k_{12} - \tfrac15$, $h_2 = k_{10} + k_{12} + k_{21} - \tfrac45$, $h_3 = k_{10}k_{21} - \tfrac{8}{100}$.

\schritt{Schritt 2 --- S-Polynome}
$S(h_2, h_3) = k_{21}h_2 - h_3 = k_{12}k_{21} + k_{21}^{2} - \tfrac45k_{21} + \tfrac{8}{100}$. Rest durch $\{h_1, h_2, h_3\}$: $k_{12}k_{21}$ ist durch $k_{12}$ teilbar: $- k_{21}h_1$ gibt $k_{21}^{2} - \tfrac45k_{21} + \tfrac15k_{21} + \tfrac{8}{100} = k_{21}^{2} - \tfrac35k_{21} + \tfrac{8}{100}$; kein Leitmonom passt: neu, $h_4$. $S(h_1, h_2) = k_{10}h_1 - k_{12}h_2 = -k_{12}^{2} - k_{12}k_{21} - \tfrac15k_{10} + \tfrac45k_{12}$; Rest: alles reduzierbar auf $0$ (nachrechnen: $-k_{12}^{2}\to$ mit $h_1$ \ldots\ ergibt am Ende $-\tfrac15h_2$-Vielfache); die übrigen Paare ebenso.

\schritt{Schritt 3 --- Reduzieren}
$h_3$ hat Leitmonom $k_{10}k_{21}$, teilbar durch $\mathrm{LM}(h_2) = k_{10}$: fliegt heraus. $h_2$ reduziert durch $h_1$: $k_{10} + k_{21} - \tfrac35$. Basis:
\[ G = \bigl\{\,k_{10} + k_{21} - \tfrac35,\quad k_{12} - \tfrac15,\quad k_{21}^{2} - \tfrac35k_{21} + \tfrac{8}{100}\,\bigr\} . \]

\schritt{Schritt 4 --- Lesen}
Leitmonome $k_{10}$, $k_{12}$, $k_{21}^{2}$: Kein Parameter ist frei (jeder kommt in einem Leitmonom vor) --- Dimension null, endlich viele Lösungen. Das letzte Element hat Grad $2$ in $k_{21}$: zwei Lösungen, $k_{21} = \tfrac15$ oder $\tfrac25$ (Aufgabe 1.3). \checkmark\ Lokal, nicht global identifizierbar; die Zahl der Lösungen ist das Produkt der Leitmonom-Grade: $1\cdot1\cdot2 = 2$.

\begin{center}
\begin{minipage}{0.88\textwidth}
\kursiv{\textbf{Die Regel \enquote{Zahl der Lösungen = Produkt der Grade der Leitmonome} gilt für lex-Basen in Dreiecksform} (jede Veränderliche genau ein Leitmonom, Potenz einer Veränderlichen) und ist ein Spezialfall des Satzes, dass die Zahl der Standardmonome die Zahl der Lösungen (mit Vielfachheit) ist, Kapitel 7. Die Programme zur Identifizierbarkeit geben genau diese Zahl aus: $1$ heißt global, $2$ oder mehr lokal, $\infty$ nicht identifizierbar.}
\end{minipage}
\end{center}
\end{aufgabe}

\begin{geschichte}
\boxtitelg{Innsbruck 1965}
\Person{Bruno Buchberger} schrieb 1965 in Innsbruck seine Dissertation bei \Person{Wolfgang Gröbner} über die Frage, wie man den Restklassenring $k[x_1, \ldots, x_n]/I$ berechnet --- eine Basis des Vektorraums, eine Multiplikationstafel. Die Antwort war das Kriterium und der Algorithmus dieses Kapitels; die Basen nannte er nach seinem Lehrer.

\begin{center}
\begin{minipage}{0.9\textwidth}
\kursiv{Die Arbeit erschien 1970 in den Aequationes Mathematicae und blieb zehn Jahre fast unbeachtet; erst mit den Computeralgebrasystemen der achtziger Jahre --- \Person{Buchberger}s eigenem in Linz, Macsyma, Reduce --- wurde sie zum Fundament eines Gebiets. \textbf{Dasselbe Verfahren war mehrfach erfunden worden}: \Person{Hironaka} 1964 für Potenzreihenringe (\enquote{Standardbasen}) im Beweis der Auflösung von Singularitäten, \Person{Shirshov} 1962 für Lie-Algebren, \Person{Bergman} 1978 für nichtkommutative Ringe. Buchbergers Fassung ist die, die man rechnen kann, und sein Kriterium ist der Teil, den die anderen nicht hatten.}
\end{minipage}
\end{center}

\schritt{F4 und F5}
\begin{center}
\begin{minipage}{0.9\textwidth}
\kursiv{\Person{Jean-Charles Faugère} ersetzte 1999 die einzelne Division durch Gauß-Elimination auf einer Matrix aller S-Polynome eines Grades (F4) und erkannte 2002 mit Signaturen vorab, welche S-Polynome zu null reduzieren (F5) --- für reguläre Folgen alle überflüssigen. \textbf{Damit fielen Probleme, die zwanzig Jahre offen waren}: cyclic-9 2001, und 2002 die algebraischen Angriffe auf das Kryptosystem HFE, Kapitel 14 --- Gröbner-Basen als Waffe gegen Verschlüsselung.}
\end{minipage}
\end{center}
\end{geschichte}

\begin{forschung}
\boxtitelp{Gröbner-Basen heute}
\begin{center}
{\small
\renewcommand{\arraystretch}{1.45}
\begin{tabular}{@{}P{3.0cm}P{5.2cm}@{}}
\toprule
\textbf{Komplexität} & Für reguläre Folgen ist der Grad der grevlex-Basis die Macaulay-Schranke $\sum(d_i - 1) + 1$, und F5 braucht polynomiell viele Operationen in dieser Schranke (Bardet--Faugère--Salvy 2015); außerhalb des regulären Falls ist wenig bewiesen.\\
\textbf{Modulare und numerische Methoden} & Rechnen modulo Primzahlen und Zusammensetzen (Arnold 2003); numerische Fortsetzung statt Algebra (Bertini, HomotopyContinuation.jl) --- Band 5 als Konkurrent von Kapitel 5.\\
\textbf{Maschinelles Lernen für Paarauswahl} & Welches S-Polynom als nächstes? Peifer, Stillman, Halpern-Leistner 2020: ein gelerntes Verfahren schlägt die klassischen Strategien --- eine Heuristik aus Daten in einem Algorithmus mit Beweis.\\
\bottomrule
\end{tabular}}
\end{center}

\schritt{Zur zweiten Zeile}
\begin{center}
\begin{minipage}{0.88\textwidth}
\kursiv{\textbf{Die Konkurrenz der numerischen algebraischen Geometrie ist ernst}: Statt eine Gröbner-Basis zu berechnen, die alle Lösungen auf einmal in exakter Arithmetik beschreibt, verfolgt man die Lösungen eines einfachen Startsystems stetig in die des gegebenen (Homotopie) und erhält alle isolierten Lösungen numerisch --- mit Fehlerschranken, in Gleitkomma, auch für Systeme, bei denen Buchberger nie enden würde. Für die Identifizierbarkeit großer Modelle ist das der heutige Standard (SIAN, 2019); die Gröbner-Basis bleibt das Werkzeug für die \kursiv{Struktur} --- Dimension, Komponenten, identifizierbare Funktionen ---, die Homotopie das für die \kursiv{Zahlen}.}
\end{minipage}
\end{center}
\end{forschung}

\begin{merksatz}
\boxtitel{Was Sie jetzt können}
\begin{center}
\renewcommand{\arraystretch}{1.5}
\begin{tabular}{@{}P{3.0cm}P{5.6cm}@{}}
\toprule
\textbf{Definition 5.1} & Gröbner-Basis: $\langle\mathrm{LM}(G)\rangle = \mathrm{LM}(I)$; reduziert\\
\textbf{Satz 5.2} & eindeutiger Rest; Zugehörigkeit; reduzierte Basis eindeutig\\
\textbf{Satz 5.4} & Buchberger-Kriterium: alle $\overline{S(g_i, g_j)}^{\,G} = 0$\\
\textbf{Satz 5.5} & der Algorithmus endet und ist korrekt\\
\textbf{Der Bruchfall} & Größe hängt von der Ordnung ab; lex kann explodieren\\
\bottomrule
\end{tabular}
\end{center}

\begin{center}
\begin{minipage}{0.88\textwidth}
\kursiv{Die lex-Basis des Leitfalls war dreieckig, und das war kein Zufall: lex-Basen eliminieren Veränderliche. Kapitel 6 macht daraus einen Satz --- und benutzt ihn, um Parametrisierungen in Gleichungen zu verwandeln und aus Differentialgleichungen die Messgrößen des Leitfalls herzuleiten.}
\end{minipage}
\end{center}
\end{merksatz}
